\section{Introduction}

Wholesale electricity markets clear based on the offers submitted by generation firms. These offers determine dispatch, market prices, producer revenues, and consumer payments. Because short-run electricity demand is often inelastic and supply can become tight during peak hours, strategic bidding can affect market outcomes. A large body of literature studies market power in electricity markets through various methods including but not limited to supply-function competition, simulation-based market-power measures, and empirical tests of unilateral market power \cite{green1992competition,borenstein1999market,wolak2003measuring}. Agent-based electricity-market models can also provide a complementary way to study strategic bidding in controlled simulated environments \cite{weidlich2008critical}, while strategic benchmark models are useful because different oligopoly assumptions can imply different market outcomes \cite{newbery2017strategic}.

Recent progress in the ability of large language models (LLMs) raises a question for future electricity markets: if market participants use artificial intelligence (AI) to support bidding or operational decisions, how can the natural-language objectives given to those systems affect market behavior? Recent literature begins to explore this. Algorithmic pricing agents have been shown to reach supracompetitive outcomes without explicit communication \cite{calvano2020artificial}, and recent work shows that LLM-based pricing agents are sensitive to prompt wording in strategically significant ways \cite{fish2024algorithmic}. LLMs are mostly directed with natural-language task descriptions \cite{brown2020language}, and generative-agent architectures increasingly combine memory, reflection, and planning components to produce sequential strategies \cite{park2023generative}. In work closely related to this study, Li, Kim, and Chen apply structured Thought Action Reflection Journal (TARJ) prompting to power dispatch and auction settings, using LLM agents that submit bids in simultaneous ascending auctions \cite{li2026behavioral}.

This paper studies whether objective wording can be quantitatively linked to LLM-agent bidding behavior in an electricity-market setting. This matters because future AI-assisted bidding systems are heavily influenced by natural-language objectives used to guide decisions. Since model reasoning is not directly observable due to the black-box nature of commonly used LLMs, we focus on measurable artifacts to analyze behavior: prompt objectives, generated rationale text, and submitted bids.

We introduce a semantic-orientation framework based on sentence BERT (SBERT) sentence embeddings \cite{reimers2019sentencebert}. Building on semantic-axis methods \cite{an2018semaxis,kozlowski2019geometry}, we construct a profit-versus-compliance axis and use it to measure both objective prompts and generated reasoning. We then test whether semantic orientation aligns with markups, prices, profits, and benchmark comparisons in a stylized electricity-market simulation. This provides a controlled way to evaluate how prompt objectives shift LLM-agent behavior and whether generated reasoning artifacts move consistently with submitted bids. This paper makes two contributions. It treats prompt objectives as controlled interventions for studying LLM bidding behavior. It also introduces an SBERT-based semantic-orientation measure that can be used both to select prompt objectives before simulation and to test whether prompt meaning, generated reasoning, and submitted bids are directionally aligned. We believe this is a concrete first step in exploring LLM-agents' possible influence on future energy markets.


The rest of the paper continues as follows. Section II describes the market simulation, LLM bidding-agent design, prompt treatments, SBERT orientation measure, and evaluation metrics. Section III describes benchmark comparisons, treatment-level outcomes, SBERT orientation regression coefficients, and TARJ alignment results. Section IV interprets the findings and methodological implications. Section V concludes with directions for richer market settings, additional models, and broader prompt designs.
